% ============================================================
% UFUU Bell Paper — Weakened Realization Claim (Alternative to T4)
% File: UFUU_Bell_T1T3_WeakenedClaim_20260502.tex
% Author: W. Jason Tuttle (ORCID: 0009-0003-3830-0551)
% Date: 2026-05-02
% Status: Analytically complete — no computational verification required
%
% Purpose: Alternative to UFUU_Bell_T4_Uniqueness_20260502.tex.
%   Rather than proving T1-T3 uniquely select H+CNOT (which requires T4),
%   this section proves the weaker but fully sufficient claim:
%   T1-T3 select the CLASS of maximally entangling gates, every member
%   of which produces a Bell state from product inputs, and H+CNOT is
%   the canonical representative of that class.
%
%   The CHSH = 2sqrt(2) result is preserved in full, because it holds
%   for ANY maximally entangling gate, not H+CNOT specifically.
%   The weakened framing loses nothing physically relevant.
%
% Relationship to T4 file:
%   Use THIS file if you prefer the cleaner, more conservative argument.
%   Use the T4 file if you want the uniqueness claim with the additional
%   minimum-depth constraint. Both are defensible; this one is simpler.
%
% To include in the Bell paper:
%   Replace Section 3.1-3.2 framing with this section, or
%   % ============================================================
% UFUU Bell Paper — Weakened Realization Claim (Alternative to T4)
% File: UFUU_Bell_T1T3_WeakenedClaim_20260502.tex
% Author: W. Jason Tuttle (ORCID: 0009-0003-3830-0551)
% Date: 2026-05-02
% Status: Analytically complete — no computational verification required
%
% Purpose: Alternative to UFUU_Bell_T4_Uniqueness_20260502.tex.
%   Rather than proving T1-T3 uniquely select H+CNOT (which requires T4),
%   this section proves the weaker but fully sufficient claim:
%   T1-T3 select the CLASS of maximally entangling gates, every member
%   of which produces a Bell state from product inputs, and H+CNOT is
%   the canonical representative of that class.
%
%   The CHSH = 2sqrt(2) result is preserved in full, because it holds
%   for ANY maximally entangling gate, not H+CNOT specifically.
%   The weakened framing loses nothing physically relevant.
%
% Relationship to T4 file:
%   Use THIS file if you prefer the cleaner, more conservative argument.
%   Use the T4 file if you want the uniqueness claim with the additional
%   minimum-depth constraint. Both are defensible; this one is simpler.
%
% To include in the Bell paper:
%   Replace Section 3.1-3.2 framing with this section, or
%   % ============================================================
% UFUU Bell Paper — Weakened Realization Claim (Alternative to T4)
% File: UFUU_Bell_T1T3_WeakenedClaim_20260502.tex
% Author: W. Jason Tuttle (ORCID: 0009-0003-3830-0551)
% Date: 2026-05-02
% Status: Analytically complete — no computational verification required
%
% Purpose: Alternative to UFUU_Bell_T4_Uniqueness_20260502.tex.
%   Rather than proving T1-T3 uniquely select H+CNOT (which requires T4),
%   this section proves the weaker but fully sufficient claim:
%   T1-T3 select the CLASS of maximally entangling gates, every member
%   of which produces a Bell state from product inputs, and H+CNOT is
%   the canonical representative of that class.
%
%   The CHSH = 2sqrt(2) result is preserved in full, because it holds
%   for ANY maximally entangling gate, not H+CNOT specifically.
%   The weakened framing loses nothing physically relevant.
%
% Relationship to T4 file:
%   Use THIS file if you prefer the cleaner, more conservative argument.
%   Use the T4 file if you want the uniqueness claim with the additional
%   minimum-depth constraint. Both are defensible; this one is simpler.
%
% To include in the Bell paper:
%   Replace Section 3.1-3.2 framing with this section, or
%   % ============================================================
% UFUU Bell Paper — Weakened Realization Claim (Alternative to T4)
% File: UFUU_Bell_T1T3_WeakenedClaim_20260502.tex
% Author: W. Jason Tuttle (ORCID: 0009-0003-3830-0551)
% Date: 2026-05-02
% Status: Analytically complete — no computational verification required
%
% Purpose: Alternative to UFUU_Bell_T4_Uniqueness_20260502.tex.
%   Rather than proving T1-T3 uniquely select H+CNOT (which requires T4),
%   this section proves the weaker but fully sufficient claim:
%   T1-T3 select the CLASS of maximally entangling gates, every member
%   of which produces a Bell state from product inputs, and H+CNOT is
%   the canonical representative of that class.
%
%   The CHSH = 2sqrt(2) result is preserved in full, because it holds
%   for ANY maximally entangling gate, not H+CNOT specifically.
%   The weakened framing loses nothing physically relevant.
%
% Relationship to T4 file:
%   Use THIS file if you prefer the cleaner, more conservative argument.
%   Use the T4 file if you want the uniqueness claim with the additional
%   minimum-depth constraint. Both are defensible; this one is simpler.
%
% To include in the Bell paper:
%   Replace Section 3.1-3.2 framing with this section, or
%   \input{UFUU_Bell_T1T3_WeakenedClaim_20260502}
% ============================================================

% ── Standalone preamble — uncomment to compile alone ────────
% \documentclass[12pt]{article}
% \usepackage{amsmath,amssymb,amsthm}
% \usepackage{physics}
% \usepackage{braket}
% \usepackage{hyperref}
% \newtheorem{theorem}{Theorem}
% \newtheorem{lemma}{Lemma}
% \newtheorem{definition}{Definition}
% \newtheorem{corollary}{Corollary}
% \newtheorem{remark}{Remark}
% \begin{document}
% ── End standalone preamble ──────────────────────────────────

\subsection{Realization of the Fold in Two-Qubit Hilbert Space}
\label{sec:fold_realization_weak}

The abstract fold operation $F$, constrained by Tuttle's Triad (T1--T3),
must be realized as a concrete quantum circuit in the two-qubit Hilbert
space $\mathcal{H} = \mathbb{C}^2 \otimes \mathbb{C}^2$. We prove that
T1--T3 jointly select the \emph{class} of maximally entangling gates ---
every member of which maps product inputs to Bell states and yields
$\langle\mathrm{CHSH}\rangle = 2\sqrt{2}$ --- and identify
$F_{\mathrm{circuit}} = \mathrm{CNOT}\cdot(H\otimes I)$ as the
canonical representative of this class.

% ── Key definition ────────────────────────────────────────────
\begin{definition}[Maximally Entangling Gate / Perfect Entangler]
\label{def:perfect_entangler_weak}
A two-qubit gate $U \in SU(4)$ is a \emph{maximally entangling gate}
(equivalently, a \emph{perfect entangler}~\cite{Zhang2003}) if there
exists a product state $\ket{\psi} = \ket{a}\otimes\ket{b}$ such that
$U\ket{\psi}$ has concurrence $\mathcal{C} = 1$, i.e., is a maximally
entangled (Bell-class) state.
\end{definition}

\begin{definition}[Bell-Class State]
A two-qubit pure state is a \emph{Bell-class state} if it is related
to one of the four Bell states
$\{\ket{\Phi^\pm}, \ket{\Psi^\pm}\}$ by local (single-qubit) unitaries.
All Bell-class states have concurrence $\mathcal{C} = 1$ and yield
$\langle\mathrm{CHSH}\rangle = 2\sqrt{2}$ under optimal measurement
settings~\cite{Cirelson1980}.
\end{definition}

% ── Main theorem (weakened) ────────────────────────────────────
\begin{theorem}[T1--T3 Select the Maximally Entangling Class]
\label{thm:weak_selection}
Any two-qubit unitary realization of the fold operation $F$ satisfying
Tuttle's Triad T1 (non-commutativity), T2 (nonlinearity), and T3
(partial irreversibility) is a maximally entangling gate. Consequently,
any such realization maps two initially independent binary distinctions
$\ket{00}$ (with appropriate single-qubit preparation) to a Bell-class
state, and the resulting CHSH correlator satisfies
$\langle\mathrm{CHSH}\rangle = 2\sqrt{2}$.
\end{theorem}

\begin{proof}
We show each triad constraint eliminates a distinct subclass of gates,
and the surviving class is precisely the maximally entangling class.

\medskip
\noindent\textbf{T2 eliminates non-entangling gates.}
T2 (nonlinearity) requires that the output of $F$ is not a linear
superposition of the outputs obtained by applying $F$ to each input
independently --- that is, the gate must generate genuine two-party
correlations absent from the inputs. In quantum circuit language, $F$
must increase entanglement from a product input. Any gate that maps
every product state to a product state (a \emph{non-entangling} gate)
trivially violates T2. Therefore $F$ must be an entangling gate.
Formally, the set of non-entangling gates has measure zero in $SU(4)$
and is precisely the set of product unitaries $A\otimes B$; all of
these are eliminated.

\medskip
\noindent\textbf{T2 further requires maximal entanglement.}
The information-maximization primitive underlying T2 is not merely
that some entanglement is generated, but that the fold selects the
phase configuration maximizing local von Neumann entropy at every
recursive step (established in Section~2 and computationally
verified for the UFUU transfer operator). An entangling gate that
produces only partial entanglement ($0 < \mathcal{C} < 1$) does
not maximize the entropy of the output state given a product input.
T2 therefore requires $\mathcal{C} = 1$: the gate must be a
maximally entangling gate (Definition~\ref{def:perfect_entangler_weak}).

\medskip
\noindent\textbf{T1 eliminates symmetric gates.}
T1 (non-commutativity) requires $F(a,b) \neq F(b,a)$ --- the fold
is directional. In the two-qubit Hilbert space, this eliminates all
gates $U$ for which swapping the two subsystems leaves $U$ invariant,
i.e., $U = \mathrm{SWAP}\cdot U\cdot\mathrm{SWAP}$. The most
prominent example is the SWAP gate itself, which maps
$\ket{a}\otimes\ket{b} \mapsto \ket{b}\otimes\ket{a}$ symmetrically.
More generally, T1 removes all gates in the symmetric subspace of
$SU(4)$. Note that SWAP has Weyl coordinates $(\pi/4, \pi/4, \pi/4)$
and concurrence $\mathcal{C} = 0$ for any product input, so it is
already eliminated by T2; T1 provides independent and complementary
elimination of the symmetric class.

\medskip
\noindent\textbf{T3 is consistent with the surviving class.}
T3 (partial irreversibility) requires that the fold saturates when
further recursive applications cannot produce distinguishable new
outputs at the prevailing UV resolution scale. Individual applications
of any unitary gate are reversible; T3 applies to the \emph{accumulated
recursion} to termination depth $D$, not to single gate applications.
Every maximally entangling gate, when iterated to termination depth,
generates a fixed-point state encoding thermodynamically irreversible
classical correlations (the classical record produced at measurement).
T3 does not further restrict the class of maximally entangling gates;
it is satisfied by all of them.

\medskip
\noindent\textbf{Conclusion.}
T2 selects the maximally entangling class; T1 removes the symmetric
subclass (independently of T2); T3 is satisfied by all surviving
gates. The class selected by T1--T3 is the class of asymmetric
maximally entangling gates --- equivalently, the non-symmetric
perfect entanglers in $SU(4)$.
\qed
\end{proof}

% ── CHSH corollary ────────────────────────────────────────────
\begin{corollary}[CHSH Violation from Any T1--T3 Realization]
\label{cor:CHSH_class}
Let $U$ be any two-qubit unitary realization of $F$ satisfying T1--T3,
and let $\ket{\psi} = U\cdot(A\otimes I)\ket{00}$ for appropriate
$A \in SU(2)$. Then $\ket{\psi}$ is a Bell-class state with
concurrence $\mathcal{C} = 1$, and under optimal measurement settings
\begin{equation}
    \langle\mathrm{CHSH}\rangle = 2\sqrt{2}.
    \label{eq:CHSH_class}
\end{equation}
This result is independent of which member of the maximally entangling
class is chosen as the realization of $F$.
\end{corollary}

\begin{proof}
By Theorem~\ref{thm:weak_selection}, $U$ is a maximally entangling
gate: there exists a product input $\ket{\psi_0}$ such that
$U\ket{\psi_0}$ has $\mathcal{C} = 1$. Any state with $\mathcal{C}=1$
is a Bell-class state (related to $\ket{\Phi^+}$ by local unitaries).
The CHSH correlator for any Bell-class state achieves the Tsirelson
bound $2\sqrt{2}$ under the measurement settings
$\{A = \sigma_z,\, A' = \sigma_x,\, B = (\sigma_z+\sigma_x)/\sqrt{2},\,
B' = (\sigma_z-\sigma_x)/\sqrt{2}\}$~\cite{Cirelson1980}.
Local unitaries on Alice's or Bob's side do not change the
entanglement structure and can be absorbed into redefined
measurement settings. Therefore $\langle\mathrm{CHSH}\rangle = 2\sqrt{2}$
for any Bell-class state, independent of which maximally entangling
gate produced it. \qed
\end{proof}

% ── Canonical representative ──────────────────────────────────
\subsubsection*{The Canonical Representative: $H + \mathrm{CNOT}$}

Theorem~\ref{thm:weak_selection} selects a class, not a unique gate.
Among the infinitely many maximally entangling gates satisfying T1--T3,
we identify $F_{\mathrm{circuit}} = \mathrm{CNOT}\cdot(H\otimes I)$
as the \emph{canonical representative} on the following grounds.

\begin{enumerate}

\item \textbf{Minimum circuit depth.}
$F_{\mathrm{circuit}}$ produces $\mathcal{C}=1$ from $\ket{00}$ using
exactly 2 elementary gates (one single-qubit gate plus one two-qubit
entangling gate) --- the minimum achievable, since at least one of each
is required by T2 and the product structure of $\ket{00}$.

\item \textbf{Minimum interaction complexity.}
In the Weyl decomposition~\cite{Zhang2003}, the CNOT class has
coordinates $(\pi/4, 0, 0)$ --- a single active interaction axis
($XX$ only) at maximal coupling strength $\pi/4$. No maximally
entangling gate has fewer active axes and remains a perfect entangler
in a single application from $\ket{00}$.

\item \textbf{Exact Bell state output.}
$F_{\mathrm{circuit}}\ket{00} = \ket{\Phi^+}$ exactly:
\begin{equation}
    \mathrm{CNOT}\cdot(H\otimes I)\ket{00}
    = \mathrm{CNOT}\cdot\frac{\ket{0}+\ket{1}}{\sqrt{2}}\otimes\ket{0}
    = \frac{\ket{00}+\ket{11}}{\sqrt{2}}
    = \ket{\Phi^+}.
    \label{eq:H_CNOT_output}
\end{equation}
No additional local unitaries are required to reach the standard Bell
state; other maximally entangling gates produce Bell-class states
that may require local rotation to reach $\ket{\Phi^+}$.

\item \textbf{Physical transparency.}
The Hadamard $H$ implements the superposition step (one binary
distinction placed in uniform superposition); CNOT implements the
correlation step (the second binary distinction conditioned on the
first). This two-step structure --- superpose then correlate ---
maps directly onto the fold's recursive decomposition: a state
decomposes into left/right sub-states (superposition), which then
interact via the fold (correlation).

\end{enumerate}

\begin{remark}[What the Canonical Choice Does and Does Not Claim]
\label{rem:canonical_caveat}
The designation of $F_{\mathrm{circuit}}$ as canonical does not imply
that other maximally entangling gates are physically excluded by
T1--T3. It implies that $F_{\mathrm{circuit}}$ is the most
parsimonious representative of the T1--T3-selected class when
additional minimality criteria (circuit depth, interaction
complexity, directness of Bell state output) are applied. The
physical result --- entanglement and CHSH violation at the Tsirelson
bound --- is a property of the class, not of the representative.
The representative is chosen for clarity of exposition and
consistency with the fold's information-maximization primitive.
\end{remark}

% ── Comparison to uniqueness claim ────────────────────────────
\subsubsection*{Relationship to the Uniqueness Result}

A stronger claim --- that T1--T3, augmented by a minimum circuit
depth constraint (T4), uniquely select the CNOT Weyl class up to
local unitaries --- is established separately in the companion
note \texttt{UFUU\_Bell\_T4\_Uniqueness\_20260502.tex}. The present
section makes no use of T4 and requires only T1--T3. The physical
conclusions are identical: entanglement from product inputs and
$\langle\mathrm{CHSH}\rangle = 2\sqrt{2}$ as structural consequences
of the fold. The difference is rhetorical precision:
Theorem~\ref{thm:weak_selection} says the fold \emph{cannot avoid}
producing Bell violations; the T4 theorem says it does so via
\emph{exactly one} gate class. Either framing supports the paper's
central claim.

% ── References ────────────────────────────────────────────────
% Add to main bibliography if not already present:
%
% \bibitem{Zhang2003}
%   Zhang J, Vala J, Sastry S, Whaley K B (2003).
%   Geometric theory of nonlocal two-qubit operations.
%   \textit{Physical Review A} \textbf{67}, 042313.
%
% \bibitem{Cirelson1980}
%   Cirel'son B S (1980).
%   Quantum generalizations of Bell's inequality.
%   \textit{Letters in Mathematical Physics} \textbf{4}, 93--100.
%
% \bibitem{Childs2003}
%   Childs A M, Haselgrove H L, Nielsen M A (2003).
%   Lower bounds on the complexity of simulating quantum gates.
%   \textit{Physical Review A} \textbf{68}, 052311.

% ── End standalone ─────────────────────────────────────────────
% \end{document}

% ============================================================

% ── Standalone preamble — uncomment to compile alone ────────
% \documentclass[12pt]{article}
% \usepackage{amsmath,amssymb,amsthm}
% \usepackage{physics}
% \usepackage{braket}
% \usepackage{hyperref}
% \newtheorem{theorem}{Theorem}
% \newtheorem{lemma}{Lemma}
% \newtheorem{definition}{Definition}
% \newtheorem{corollary}{Corollary}
% \newtheorem{remark}{Remark}
% \begin{document}
% ── End standalone preamble ──────────────────────────────────

\subsection{Realization of the Fold in Two-Qubit Hilbert Space}
\label{sec:fold_realization_weak}

The abstract fold operation $F$, constrained by Tuttle's Triad (T1--T3),
must be realized as a concrete quantum circuit in the two-qubit Hilbert
space $\mathcal{H} = \mathbb{C}^2 \otimes \mathbb{C}^2$. We prove that
T1--T3 jointly select the \emph{class} of maximally entangling gates ---
every member of which maps product inputs to Bell states and yields
$\langle\mathrm{CHSH}\rangle = 2\sqrt{2}$ --- and identify
$F_{\mathrm{circuit}} = \mathrm{CNOT}\cdot(H\otimes I)$ as the
canonical representative of this class.

% ── Key definition ────────────────────────────────────────────
\begin{definition}[Maximally Entangling Gate / Perfect Entangler]
\label{def:perfect_entangler_weak}
A two-qubit gate $U \in SU(4)$ is a \emph{maximally entangling gate}
(equivalently, a \emph{perfect entangler}~\cite{Zhang2003}) if there
exists a product state $\ket{\psi} = \ket{a}\otimes\ket{b}$ such that
$U\ket{\psi}$ has concurrence $\mathcal{C} = 1$, i.e., is a maximally
entangled (Bell-class) state.
\end{definition}

\begin{definition}[Bell-Class State]
A two-qubit pure state is a \emph{Bell-class state} if it is related
to one of the four Bell states
$\{\ket{\Phi^\pm}, \ket{\Psi^\pm}\}$ by local (single-qubit) unitaries.
All Bell-class states have concurrence $\mathcal{C} = 1$ and yield
$\langle\mathrm{CHSH}\rangle = 2\sqrt{2}$ under optimal measurement
settings~\cite{Cirelson1980}.
\end{definition}

% ── Main theorem (weakened) ────────────────────────────────────
\begin{theorem}[T1--T3 Select the Maximally Entangling Class]
\label{thm:weak_selection}
Any two-qubit unitary realization of the fold operation $F$ satisfying
Tuttle's Triad T1 (non-commutativity), T2 (nonlinearity), and T3
(partial irreversibility) is a maximally entangling gate. Consequently,
any such realization maps two initially independent binary distinctions
$\ket{00}$ (with appropriate single-qubit preparation) to a Bell-class
state, and the resulting CHSH correlator satisfies
$\langle\mathrm{CHSH}\rangle = 2\sqrt{2}$.
\end{theorem}

\begin{proof}
We show each triad constraint eliminates a distinct subclass of gates,
and the surviving class is precisely the maximally entangling class.

\medskip
\noindent\textbf{T2 eliminates non-entangling gates.}
T2 (nonlinearity) requires that the output of $F$ is not a linear
superposition of the outputs obtained by applying $F$ to each input
independently --- that is, the gate must generate genuine two-party
correlations absent from the inputs. In quantum circuit language, $F$
must increase entanglement from a product input. Any gate that maps
every product state to a product state (a \emph{non-entangling} gate)
trivially violates T2. Therefore $F$ must be an entangling gate.
Formally, the set of non-entangling gates has measure zero in $SU(4)$
and is precisely the set of product unitaries $A\otimes B$; all of
these are eliminated.

\medskip
\noindent\textbf{T2 further requires maximal entanglement.}
The information-maximization primitive underlying T2 is not merely
that some entanglement is generated, but that the fold selects the
phase configuration maximizing local von Neumann entropy at every
recursive step (established in Section~2 and computationally
verified for the UFUU transfer operator). An entangling gate that
produces only partial entanglement ($0 < \mathcal{C} < 1$) does
not maximize the entropy of the output state given a product input.
T2 therefore requires $\mathcal{C} = 1$: the gate must be a
maximally entangling gate (Definition~\ref{def:perfect_entangler_weak}).

\medskip
\noindent\textbf{T1 eliminates symmetric gates.}
T1 (non-commutativity) requires $F(a,b) \neq F(b,a)$ --- the fold
is directional. In the two-qubit Hilbert space, this eliminates all
gates $U$ for which swapping the two subsystems leaves $U$ invariant,
i.e., $U = \mathrm{SWAP}\cdot U\cdot\mathrm{SWAP}$. The most
prominent example is the SWAP gate itself, which maps
$\ket{a}\otimes\ket{b} \mapsto \ket{b}\otimes\ket{a}$ symmetrically.
More generally, T1 removes all gates in the symmetric subspace of
$SU(4)$. Note that SWAP has Weyl coordinates $(\pi/4, \pi/4, \pi/4)$
and concurrence $\mathcal{C} = 0$ for any product input, so it is
already eliminated by T2; T1 provides independent and complementary
elimination of the symmetric class.

\medskip
\noindent\textbf{T3 is consistent with the surviving class.}
T3 (partial irreversibility) requires that the fold saturates when
further recursive applications cannot produce distinguishable new
outputs at the prevailing UV resolution scale. Individual applications
of any unitary gate are reversible; T3 applies to the \emph{accumulated
recursion} to termination depth $D$, not to single gate applications.
Every maximally entangling gate, when iterated to termination depth,
generates a fixed-point state encoding thermodynamically irreversible
classical correlations (the classical record produced at measurement).
T3 does not further restrict the class of maximally entangling gates;
it is satisfied by all of them.

\medskip
\noindent\textbf{Conclusion.}
T2 selects the maximally entangling class; T1 removes the symmetric
subclass (independently of T2); T3 is satisfied by all surviving
gates. The class selected by T1--T3 is the class of asymmetric
maximally entangling gates --- equivalently, the non-symmetric
perfect entanglers in $SU(4)$.
\qed
\end{proof}

% ── CHSH corollary ────────────────────────────────────────────
\begin{corollary}[CHSH Violation from Any T1--T3 Realization]
\label{cor:CHSH_class}
Let $U$ be any two-qubit unitary realization of $F$ satisfying T1--T3,
and let $\ket{\psi} = U\cdot(A\otimes I)\ket{00}$ for appropriate
$A \in SU(2)$. Then $\ket{\psi}$ is a Bell-class state with
concurrence $\mathcal{C} = 1$, and under optimal measurement settings
\begin{equation}
    \langle\mathrm{CHSH}\rangle = 2\sqrt{2}.
    \label{eq:CHSH_class}
\end{equation}
This result is independent of which member of the maximally entangling
class is chosen as the realization of $F$.
\end{corollary}

\begin{proof}
By Theorem~\ref{thm:weak_selection}, $U$ is a maximally entangling
gate: there exists a product input $\ket{\psi_0}$ such that
$U\ket{\psi_0}$ has $\mathcal{C} = 1$. Any state with $\mathcal{C}=1$
is a Bell-class state (related to $\ket{\Phi^+}$ by local unitaries).
The CHSH correlator for any Bell-class state achieves the Tsirelson
bound $2\sqrt{2}$ under the measurement settings
$\{A = \sigma_z,\, A' = \sigma_x,\, B = (\sigma_z+\sigma_x)/\sqrt{2},\,
B' = (\sigma_z-\sigma_x)/\sqrt{2}\}$~\cite{Cirelson1980}.
Local unitaries on Alice's or Bob's side do not change the
entanglement structure and can be absorbed into redefined
measurement settings. Therefore $\langle\mathrm{CHSH}\rangle = 2\sqrt{2}$
for any Bell-class state, independent of which maximally entangling
gate produced it. \qed
\end{proof}

% ── Canonical representative ──────────────────────────────────
\subsubsection*{The Canonical Representative: $H + \mathrm{CNOT}$}

Theorem~\ref{thm:weak_selection} selects a class, not a unique gate.
Among the infinitely many maximally entangling gates satisfying T1--T3,
we identify $F_{\mathrm{circuit}} = \mathrm{CNOT}\cdot(H\otimes I)$
as the \emph{canonical representative} on the following grounds.

\begin{enumerate}

\item \textbf{Minimum circuit depth.}
$F_{\mathrm{circuit}}$ produces $\mathcal{C}=1$ from $\ket{00}$ using
exactly 2 elementary gates (one single-qubit gate plus one two-qubit
entangling gate) --- the minimum achievable, since at least one of each
is required by T2 and the product structure of $\ket{00}$.

\item \textbf{Minimum interaction complexity.}
In the Weyl decomposition~\cite{Zhang2003}, the CNOT class has
coordinates $(\pi/4, 0, 0)$ --- a single active interaction axis
($XX$ only) at maximal coupling strength $\pi/4$. No maximally
entangling gate has fewer active axes and remains a perfect entangler
in a single application from $\ket{00}$.

\item \textbf{Exact Bell state output.}
$F_{\mathrm{circuit}}\ket{00} = \ket{\Phi^+}$ exactly:
\begin{equation}
    \mathrm{CNOT}\cdot(H\otimes I)\ket{00}
    = \mathrm{CNOT}\cdot\frac{\ket{0}+\ket{1}}{\sqrt{2}}\otimes\ket{0}
    = \frac{\ket{00}+\ket{11}}{\sqrt{2}}
    = \ket{\Phi^+}.
    \label{eq:H_CNOT_output}
\end{equation}
No additional local unitaries are required to reach the standard Bell
state; other maximally entangling gates produce Bell-class states
that may require local rotation to reach $\ket{\Phi^+}$.

\item \textbf{Physical transparency.}
The Hadamard $H$ implements the superposition step (one binary
distinction placed in uniform superposition); CNOT implements the
correlation step (the second binary distinction conditioned on the
first). This two-step structure --- superpose then correlate ---
maps directly onto the fold's recursive decomposition: a state
decomposes into left/right sub-states (superposition), which then
interact via the fold (correlation).

\end{enumerate}

\begin{remark}[What the Canonical Choice Does and Does Not Claim]
\label{rem:canonical_caveat}
The designation of $F_{\mathrm{circuit}}$ as canonical does not imply
that other maximally entangling gates are physically excluded by
T1--T3. It implies that $F_{\mathrm{circuit}}$ is the most
parsimonious representative of the T1--T3-selected class when
additional minimality criteria (circuit depth, interaction
complexity, directness of Bell state output) are applied. The
physical result --- entanglement and CHSH violation at the Tsirelson
bound --- is a property of the class, not of the representative.
The representative is chosen for clarity of exposition and
consistency with the fold's information-maximization primitive.
\end{remark}

% ── Comparison to uniqueness claim ────────────────────────────
\subsubsection*{Relationship to the Uniqueness Result}

A stronger claim --- that T1--T3, augmented by a minimum circuit
depth constraint (T4), uniquely select the CNOT Weyl class up to
local unitaries --- is established separately in the companion
note \texttt{UFUU\_Bell\_T4\_Uniqueness\_20260502.tex}. The present
section makes no use of T4 and requires only T1--T3. The physical
conclusions are identical: entanglement from product inputs and
$\langle\mathrm{CHSH}\rangle = 2\sqrt{2}$ as structural consequences
of the fold. The difference is rhetorical precision:
Theorem~\ref{thm:weak_selection} says the fold \emph{cannot avoid}
producing Bell violations; the T4 theorem says it does so via
\emph{exactly one} gate class. Either framing supports the paper's
central claim.

% ── References ────────────────────────────────────────────────
% Add to main bibliography if not already present:
%
% \bibitem{Zhang2003}
%   Zhang J, Vala J, Sastry S, Whaley K B (2003).
%   Geometric theory of nonlocal two-qubit operations.
%   \textit{Physical Review A} \textbf{67}, 042313.
%
% \bibitem{Cirelson1980}
%   Cirel'son B S (1980).
%   Quantum generalizations of Bell's inequality.
%   \textit{Letters in Mathematical Physics} \textbf{4}, 93--100.
%
% \bibitem{Childs2003}
%   Childs A M, Haselgrove H L, Nielsen M A (2003).
%   Lower bounds on the complexity of simulating quantum gates.
%   \textit{Physical Review A} \textbf{68}, 052311.

% ── End standalone ─────────────────────────────────────────────
% \end{document}

% ============================================================

% ── Standalone preamble — uncomment to compile alone ────────
% \documentclass[12pt]{article}
% \usepackage{amsmath,amssymb,amsthm}
% \usepackage{physics}
% \usepackage{braket}
% \usepackage{hyperref}
% \newtheorem{theorem}{Theorem}
% \newtheorem{lemma}{Lemma}
% \newtheorem{definition}{Definition}
% \newtheorem{corollary}{Corollary}
% \newtheorem{remark}{Remark}
% \begin{document}
% ── End standalone preamble ──────────────────────────────────

\subsection{Realization of the Fold in Two-Qubit Hilbert Space}
\label{sec:fold_realization_weak}

The abstract fold operation $F$, constrained by Tuttle's Triad (T1--T3),
must be realized as a concrete quantum circuit in the two-qubit Hilbert
space $\mathcal{H} = \mathbb{C}^2 \otimes \mathbb{C}^2$. We prove that
T1--T3 jointly select the \emph{class} of maximally entangling gates ---
every member of which maps product inputs to Bell states and yields
$\langle\mathrm{CHSH}\rangle = 2\sqrt{2}$ --- and identify
$F_{\mathrm{circuit}} = \mathrm{CNOT}\cdot(H\otimes I)$ as the
canonical representative of this class.

% ── Key definition ────────────────────────────────────────────
\begin{definition}[Maximally Entangling Gate / Perfect Entangler]
\label{def:perfect_entangler_weak}
A two-qubit gate $U \in SU(4)$ is a \emph{maximally entangling gate}
(equivalently, a \emph{perfect entangler}~\cite{Zhang2003}) if there
exists a product state $\ket{\psi} = \ket{a}\otimes\ket{b}$ such that
$U\ket{\psi}$ has concurrence $\mathcal{C} = 1$, i.e., is a maximally
entangled (Bell-class) state.
\end{definition}

\begin{definition}[Bell-Class State]
A two-qubit pure state is a \emph{Bell-class state} if it is related
to one of the four Bell states
$\{\ket{\Phi^\pm}, \ket{\Psi^\pm}\}$ by local (single-qubit) unitaries.
All Bell-class states have concurrence $\mathcal{C} = 1$ and yield
$\langle\mathrm{CHSH}\rangle = 2\sqrt{2}$ under optimal measurement
settings~\cite{Cirelson1980}.
\end{definition}

% ── Main theorem (weakened) ────────────────────────────────────
\begin{theorem}[T1--T3 Select the Maximally Entangling Class]
\label{thm:weak_selection}
Any two-qubit unitary realization of the fold operation $F$ satisfying
Tuttle's Triad T1 (non-commutativity), T2 (nonlinearity), and T3
(partial irreversibility) is a maximally entangling gate. Consequently,
any such realization maps two initially independent binary distinctions
$\ket{00}$ (with appropriate single-qubit preparation) to a Bell-class
state, and the resulting CHSH correlator satisfies
$\langle\mathrm{CHSH}\rangle = 2\sqrt{2}$.
\end{theorem}

\begin{proof}
We show each triad constraint eliminates a distinct subclass of gates,
and the surviving class is precisely the maximally entangling class.

\medskip
\noindent\textbf{T2 eliminates non-entangling gates.}
T2 (nonlinearity) requires that the output of $F$ is not a linear
superposition of the outputs obtained by applying $F$ to each input
independently --- that is, the gate must generate genuine two-party
correlations absent from the inputs. In quantum circuit language, $F$
must increase entanglement from a product input. Any gate that maps
every product state to a product state (a \emph{non-entangling} gate)
trivially violates T2. Therefore $F$ must be an entangling gate.
Formally, the set of non-entangling gates has measure zero in $SU(4)$
and is precisely the set of product unitaries $A\otimes B$; all of
these are eliminated.

\medskip
\noindent\textbf{T2 further requires maximal entanglement.}
The information-maximization primitive underlying T2 is not merely
that some entanglement is generated, but that the fold selects the
phase configuration maximizing local von Neumann entropy at every
recursive step (established in Section~2 and computationally
verified for the UFUU transfer operator). An entangling gate that
produces only partial entanglement ($0 < \mathcal{C} < 1$) does
not maximize the entropy of the output state given a product input.
T2 therefore requires $\mathcal{C} = 1$: the gate must be a
maximally entangling gate (Definition~\ref{def:perfect_entangler_weak}).

\medskip
\noindent\textbf{T1 eliminates symmetric gates.}
T1 (non-commutativity) requires $F(a,b) \neq F(b,a)$ --- the fold
is directional. In the two-qubit Hilbert space, this eliminates all
gates $U$ for which swapping the two subsystems leaves $U$ invariant,
i.e., $U = \mathrm{SWAP}\cdot U\cdot\mathrm{SWAP}$. The most
prominent example is the SWAP gate itself, which maps
$\ket{a}\otimes\ket{b} \mapsto \ket{b}\otimes\ket{a}$ symmetrically.
More generally, T1 removes all gates in the symmetric subspace of
$SU(4)$. Note that SWAP has Weyl coordinates $(\pi/4, \pi/4, \pi/4)$
and concurrence $\mathcal{C} = 0$ for any product input, so it is
already eliminated by T2; T1 provides independent and complementary
elimination of the symmetric class.

\medskip
\noindent\textbf{T3 is consistent with the surviving class.}
T3 (partial irreversibility) requires that the fold saturates when
further recursive applications cannot produce distinguishable new
outputs at the prevailing UV resolution scale. Individual applications
of any unitary gate are reversible; T3 applies to the \emph{accumulated
recursion} to termination depth $D$, not to single gate applications.
Every maximally entangling gate, when iterated to termination depth,
generates a fixed-point state encoding thermodynamically irreversible
classical correlations (the classical record produced at measurement).
T3 does not further restrict the class of maximally entangling gates;
it is satisfied by all of them.

\medskip
\noindent\textbf{Conclusion.}
T2 selects the maximally entangling class; T1 removes the symmetric
subclass (independently of T2); T3 is satisfied by all surviving
gates. The class selected by T1--T3 is the class of asymmetric
maximally entangling gates --- equivalently, the non-symmetric
perfect entanglers in $SU(4)$.
\qed
\end{proof}

% ── CHSH corollary ────────────────────────────────────────────
\begin{corollary}[CHSH Violation from Any T1--T3 Realization]
\label{cor:CHSH_class}
Let $U$ be any two-qubit unitary realization of $F$ satisfying T1--T3,
and let $\ket{\psi} = U\cdot(A\otimes I)\ket{00}$ for appropriate
$A \in SU(2)$. Then $\ket{\psi}$ is a Bell-class state with
concurrence $\mathcal{C} = 1$, and under optimal measurement settings
\begin{equation}
    \langle\mathrm{CHSH}\rangle = 2\sqrt{2}.
    \label{eq:CHSH_class}
\end{equation}
This result is independent of which member of the maximally entangling
class is chosen as the realization of $F$.
\end{corollary}

\begin{proof}
By Theorem~\ref{thm:weak_selection}, $U$ is a maximally entangling
gate: there exists a product input $\ket{\psi_0}$ such that
$U\ket{\psi_0}$ has $\mathcal{C} = 1$. Any state with $\mathcal{C}=1$
is a Bell-class state (related to $\ket{\Phi^+}$ by local unitaries).
The CHSH correlator for any Bell-class state achieves the Tsirelson
bound $2\sqrt{2}$ under the measurement settings
$\{A = \sigma_z,\, A' = \sigma_x,\, B = (\sigma_z+\sigma_x)/\sqrt{2},\,
B' = (\sigma_z-\sigma_x)/\sqrt{2}\}$~\cite{Cirelson1980}.
Local unitaries on Alice's or Bob's side do not change the
entanglement structure and can be absorbed into redefined
measurement settings. Therefore $\langle\mathrm{CHSH}\rangle = 2\sqrt{2}$
for any Bell-class state, independent of which maximally entangling
gate produced it. \qed
\end{proof}

% ── Canonical representative ──────────────────────────────────
\subsubsection*{The Canonical Representative: $H + \mathrm{CNOT}$}

Theorem~\ref{thm:weak_selection} selects a class, not a unique gate.
Among the infinitely many maximally entangling gates satisfying T1--T3,
we identify $F_{\mathrm{circuit}} = \mathrm{CNOT}\cdot(H\otimes I)$
as the \emph{canonical representative} on the following grounds.

\begin{enumerate}

\item \textbf{Minimum circuit depth.}
$F_{\mathrm{circuit}}$ produces $\mathcal{C}=1$ from $\ket{00}$ using
exactly 2 elementary gates (one single-qubit gate plus one two-qubit
entangling gate) --- the minimum achievable, since at least one of each
is required by T2 and the product structure of $\ket{00}$.

\item \textbf{Minimum interaction complexity.}
In the Weyl decomposition~\cite{Zhang2003}, the CNOT class has
coordinates $(\pi/4, 0, 0)$ --- a single active interaction axis
($XX$ only) at maximal coupling strength $\pi/4$. No maximally
entangling gate has fewer active axes and remains a perfect entangler
in a single application from $\ket{00}$.

\item \textbf{Exact Bell state output.}
$F_{\mathrm{circuit}}\ket{00} = \ket{\Phi^+}$ exactly:
\begin{equation}
    \mathrm{CNOT}\cdot(H\otimes I)\ket{00}
    = \mathrm{CNOT}\cdot\frac{\ket{0}+\ket{1}}{\sqrt{2}}\otimes\ket{0}
    = \frac{\ket{00}+\ket{11}}{\sqrt{2}}
    = \ket{\Phi^+}.
    \label{eq:H_CNOT_output}
\end{equation}
No additional local unitaries are required to reach the standard Bell
state; other maximally entangling gates produce Bell-class states
that may require local rotation to reach $\ket{\Phi^+}$.

\item \textbf{Physical transparency.}
The Hadamard $H$ implements the superposition step (one binary
distinction placed in uniform superposition); CNOT implements the
correlation step (the second binary distinction conditioned on the
first). This two-step structure --- superpose then correlate ---
maps directly onto the fold's recursive decomposition: a state
decomposes into left/right sub-states (superposition), which then
interact via the fold (correlation).

\end{enumerate}

\begin{remark}[What the Canonical Choice Does and Does Not Claim]
\label{rem:canonical_caveat}
The designation of $F_{\mathrm{circuit}}$ as canonical does not imply
that other maximally entangling gates are physically excluded by
T1--T3. It implies that $F_{\mathrm{circuit}}$ is the most
parsimonious representative of the T1--T3-selected class when
additional minimality criteria (circuit depth, interaction
complexity, directness of Bell state output) are applied. The
physical result --- entanglement and CHSH violation at the Tsirelson
bound --- is a property of the class, not of the representative.
The representative is chosen for clarity of exposition and
consistency with the fold's information-maximization primitive.
\end{remark}

% ── Comparison to uniqueness claim ────────────────────────────
\subsubsection*{Relationship to the Uniqueness Result}

A stronger claim --- that T1--T3, augmented by a minimum circuit
depth constraint (T4), uniquely select the CNOT Weyl class up to
local unitaries --- is established separately in the companion
note \texttt{UFUU\_Bell\_T4\_Uniqueness\_20260502.tex}. The present
section makes no use of T4 and requires only T1--T3. The physical
conclusions are identical: entanglement from product inputs and
$\langle\mathrm{CHSH}\rangle = 2\sqrt{2}$ as structural consequences
of the fold. The difference is rhetorical precision:
Theorem~\ref{thm:weak_selection} says the fold \emph{cannot avoid}
producing Bell violations; the T4 theorem says it does so via
\emph{exactly one} gate class. Either framing supports the paper's
central claim.

% ── References ────────────────────────────────────────────────
% Add to main bibliography if not already present:
%
% \bibitem{Zhang2003}
%   Zhang J, Vala J, Sastry S, Whaley K B (2003).
%   Geometric theory of nonlocal two-qubit operations.
%   \textit{Physical Review A} \textbf{67}, 042313.
%
% \bibitem{Cirelson1980}
%   Cirel'son B S (1980).
%   Quantum generalizations of Bell's inequality.
%   \textit{Letters in Mathematical Physics} \textbf{4}, 93--100.
%
% \bibitem{Childs2003}
%   Childs A M, Haselgrove H L, Nielsen M A (2003).
%   Lower bounds on the complexity of simulating quantum gates.
%   \textit{Physical Review A} \textbf{68}, 052311.

% ── End standalone ─────────────────────────────────────────────
% \end{document}

% ============================================================

% ── Standalone preamble — uncomment to compile alone ────────
% \documentclass[12pt]{article}
% \usepackage{amsmath,amssymb,amsthm}
% \usepackage{physics}
% \usepackage{braket}
% \usepackage{hyperref}
% \newtheorem{theorem}{Theorem}
% \newtheorem{lemma}{Lemma}
% \newtheorem{definition}{Definition}
% \newtheorem{corollary}{Corollary}
% \newtheorem{remark}{Remark}
% \begin{document}
% ── End standalone preamble ──────────────────────────────────

\subsection{Realization of the Fold in Two-Qubit Hilbert Space}
\label{sec:fold_realization_weak}

The abstract fold operation $F$, constrained by Tuttle's Triad (T1--T3),
must be realized as a concrete quantum circuit in the two-qubit Hilbert
space $\mathcal{H} = \mathbb{C}^2 \otimes \mathbb{C}^2$. We prove that
T1--T3 jointly select the \emph{class} of maximally entangling gates ---
every member of which maps product inputs to Bell states and yields
$\langle\mathrm{CHSH}\rangle = 2\sqrt{2}$ --- and identify
$F_{\mathrm{circuit}} = \mathrm{CNOT}\cdot(H\otimes I)$ as the
canonical representative of this class.

% ── Key definition ────────────────────────────────────────────
\begin{definition}[Maximally Entangling Gate / Perfect Entangler]
\label{def:perfect_entangler_weak}
A two-qubit gate $U \in SU(4)$ is a \emph{maximally entangling gate}
(equivalently, a \emph{perfect entangler}~\cite{Zhang2003}) if there
exists a product state $\ket{\psi} = \ket{a}\otimes\ket{b}$ such that
$U\ket{\psi}$ has concurrence $\mathcal{C} = 1$, i.e., is a maximally
entangled (Bell-class) state.
\end{definition}

\begin{definition}[Bell-Class State]
A two-qubit pure state is a \emph{Bell-class state} if it is related
to one of the four Bell states
$\{\ket{\Phi^\pm}, \ket{\Psi^\pm}\}$ by local (single-qubit) unitaries.
All Bell-class states have concurrence $\mathcal{C} = 1$ and yield
$\langle\mathrm{CHSH}\rangle = 2\sqrt{2}$ under optimal measurement
settings~\cite{Cirelson1980}.
\end{definition}

% ── Main theorem (weakened) ────────────────────────────────────
\begin{theorem}[T1--T3 Select the Maximally Entangling Class]
\label{thm:weak_selection}
Any two-qubit unitary realization of the fold operation $F$ satisfying
Tuttle's Triad T1 (non-commutativity), T2 (nonlinearity), and T3
(partial irreversibility) is a maximally entangling gate. Consequently,
any such realization maps two initially independent binary distinctions
$\ket{00}$ (with appropriate single-qubit preparation) to a Bell-class
state, and the resulting CHSH correlator satisfies
$\langle\mathrm{CHSH}\rangle = 2\sqrt{2}$.
\end{theorem}

\begin{proof}
We show each triad constraint eliminates a distinct subclass of gates,
and the surviving class is precisely the maximally entangling class.

\medskip
\noindent\textbf{T2 eliminates non-entangling gates.}
T2 (nonlinearity) requires that the output of $F$ is not a linear
superposition of the outputs obtained by applying $F$ to each input
independently --- that is, the gate must generate genuine two-party
correlations absent from the inputs. In quantum circuit language, $F$
must increase entanglement from a product input. Any gate that maps
every product state to a product state (a \emph{non-entangling} gate)
trivially violates T2. Therefore $F$ must be an entangling gate.
Formally, the set of non-entangling gates has measure zero in $SU(4)$
and is precisely the set of product unitaries $A\otimes B$; all of
these are eliminated.

\medskip
\noindent\textbf{T2 further requires maximal entanglement.}
The information-maximization primitive underlying T2 is not merely
that some entanglement is generated, but that the fold selects the
phase configuration maximizing local von Neumann entropy at every
recursive step (established in Section~2 and computationally
verified for the UFUU transfer operator). An entangling gate that
produces only partial entanglement ($0 < \mathcal{C} < 1$) does
not maximize the entropy of the output state given a product input.
T2 therefore requires $\mathcal{C} = 1$: the gate must be a
maximally entangling gate (Definition~\ref{def:perfect_entangler_weak}).

\medskip
\noindent\textbf{T1 eliminates symmetric gates.}
T1 (non-commutativity) requires $F(a,b) \neq F(b,a)$ --- the fold
is directional. In the two-qubit Hilbert space, this eliminates all
gates $U$ for which swapping the two subsystems leaves $U$ invariant,
i.e., $U = \mathrm{SWAP}\cdot U\cdot\mathrm{SWAP}$. The most
prominent example is the SWAP gate itself, which maps
$\ket{a}\otimes\ket{b} \mapsto \ket{b}\otimes\ket{a}$ symmetrically.
More generally, T1 removes all gates in the symmetric subspace of
$SU(4)$. Note that SWAP has Weyl coordinates $(\pi/4, \pi/4, \pi/4)$
and concurrence $\mathcal{C} = 0$ for any product input, so it is
already eliminated by T2; T1 provides independent and complementary
elimination of the symmetric class.

\medskip
\noindent\textbf{T3 is consistent with the surviving class.}
T3 (partial irreversibility) requires that the fold saturates when
further recursive applications cannot produce distinguishable new
outputs at the prevailing UV resolution scale. Individual applications
of any unitary gate are reversible; T3 applies to the \emph{accumulated
recursion} to termination depth $D$, not to single gate applications.
Every maximally entangling gate, when iterated to termination depth,
generates a fixed-point state encoding thermodynamically irreversible
classical correlations (the classical record produced at measurement).
T3 does not further restrict the class of maximally entangling gates;
it is satisfied by all of them.

\medskip
\noindent\textbf{Conclusion.}
T2 selects the maximally entangling class; T1 removes the symmetric
subclass (independently of T2); T3 is satisfied by all surviving
gates. The class selected by T1--T3 is the class of asymmetric
maximally entangling gates --- equivalently, the non-symmetric
perfect entanglers in $SU(4)$.
\qed
\end{proof}

% ── CHSH corollary ────────────────────────────────────────────
\begin{corollary}[CHSH Violation from Any T1--T3 Realization]
\label{cor:CHSH_class}
Let $U$ be any two-qubit unitary realization of $F$ satisfying T1--T3,
and let $\ket{\psi} = U\cdot(A\otimes I)\ket{00}$ for appropriate
$A \in SU(2)$. Then $\ket{\psi}$ is a Bell-class state with
concurrence $\mathcal{C} = 1$, and under optimal measurement settings
\begin{equation}
    \langle\mathrm{CHSH}\rangle = 2\sqrt{2}.
    \label{eq:CHSH_class}
\end{equation}
This result is independent of which member of the maximally entangling
class is chosen as the realization of $F$.
\end{corollary}

\begin{proof}
By Theorem~\ref{thm:weak_selection}, $U$ is a maximally entangling
gate: there exists a product input $\ket{\psi_0}$ such that
$U\ket{\psi_0}$ has $\mathcal{C} = 1$. Any state with $\mathcal{C}=1$
is a Bell-class state (related to $\ket{\Phi^+}$ by local unitaries).
The CHSH correlator for any Bell-class state achieves the Tsirelson
bound $2\sqrt{2}$ under the measurement settings
$\{A = \sigma_z,\, A' = \sigma_x,\, B = (\sigma_z+\sigma_x)/\sqrt{2},\,
B' = (\sigma_z-\sigma_x)/\sqrt{2}\}$~\cite{Cirelson1980}.
Local unitaries on Alice's or Bob's side do not change the
entanglement structure and can be absorbed into redefined
measurement settings. Therefore $\langle\mathrm{CHSH}\rangle = 2\sqrt{2}$
for any Bell-class state, independent of which maximally entangling
gate produced it. \qed
\end{proof}

% ── Canonical representative ──────────────────────────────────
\subsubsection*{The Canonical Representative: $H + \mathrm{CNOT}$}

Theorem~\ref{thm:weak_selection} selects a class, not a unique gate.
Among the infinitely many maximally entangling gates satisfying T1--T3,
we identify $F_{\mathrm{circuit}} = \mathrm{CNOT}\cdot(H\otimes I)$
as the \emph{canonical representative} on the following grounds.

\begin{enumerate}

\item \textbf{Minimum circuit depth.}
$F_{\mathrm{circuit}}$ produces $\mathcal{C}=1$ from $\ket{00}$ using
exactly 2 elementary gates (one single-qubit gate plus one two-qubit
entangling gate) --- the minimum achievable, since at least one of each
is required by T2 and the product structure of $\ket{00}$.

\item \textbf{Minimum interaction complexity.}
In the Weyl decomposition~\cite{Zhang2003}, the CNOT class has
coordinates $(\pi/4, 0, 0)$ --- a single active interaction axis
($XX$ only) at maximal coupling strength $\pi/4$. No maximally
entangling gate has fewer active axes and remains a perfect entangler
in a single application from $\ket{00}$.

\item \textbf{Exact Bell state output.}
$F_{\mathrm{circuit}}\ket{00} = \ket{\Phi^+}$ exactly:
\begin{equation}
    \mathrm{CNOT}\cdot(H\otimes I)\ket{00}
    = \mathrm{CNOT}\cdot\frac{\ket{0}+\ket{1}}{\sqrt{2}}\otimes\ket{0}
    = \frac{\ket{00}+\ket{11}}{\sqrt{2}}
    = \ket{\Phi^+}.
    \label{eq:H_CNOT_output}
\end{equation}
No additional local unitaries are required to reach the standard Bell
state; other maximally entangling gates produce Bell-class states
that may require local rotation to reach $\ket{\Phi^+}$.

\item \textbf{Physical transparency.}
The Hadamard $H$ implements the superposition step (one binary
distinction placed in uniform superposition); CNOT implements the
correlation step (the second binary distinction conditioned on the
first). This two-step structure --- superpose then correlate ---
maps directly onto the fold's recursive decomposition: a state
decomposes into left/right sub-states (superposition), which then
interact via the fold (correlation).

\end{enumerate}

\begin{remark}[What the Canonical Choice Does and Does Not Claim]
\label{rem:canonical_caveat}
The designation of $F_{\mathrm{circuit}}$ as canonical does not imply
that other maximally entangling gates are physically excluded by
T1--T3. It implies that $F_{\mathrm{circuit}}$ is the most
parsimonious representative of the T1--T3-selected class when
additional minimality criteria (circuit depth, interaction
complexity, directness of Bell state output) are applied. The
physical result --- entanglement and CHSH violation at the Tsirelson
bound --- is a property of the class, not of the representative.
The representative is chosen for clarity of exposition and
consistency with the fold's information-maximization primitive.
\end{remark}

% ── Comparison to uniqueness claim ────────────────────────────
\subsubsection*{Relationship to the Uniqueness Result}

A stronger claim --- that T1--T3, augmented by a minimum circuit
depth constraint (T4), uniquely select the CNOT Weyl class up to
local unitaries --- is established separately in the companion
note \texttt{UFUU\_Bell\_T4\_Uniqueness\_20260502.tex}. The present
section makes no use of T4 and requires only T1--T3. The physical
conclusions are identical: entanglement from product inputs and
$\langle\mathrm{CHSH}\rangle = 2\sqrt{2}$ as structural consequences
of the fold. The difference is rhetorical precision:
Theorem~\ref{thm:weak_selection} says the fold \emph{cannot avoid}
producing Bell violations; the T4 theorem says it does so via
\emph{exactly one} gate class. Either framing supports the paper's
central claim.

% ── References ────────────────────────────────────────────────
% Add to main bibliography if not already present:
%
% \bibitem{Zhang2003}
%   Zhang J, Vala J, Sastry S, Whaley K B (2003).
%   Geometric theory of nonlocal two-qubit operations.
%   \textit{Physical Review A} \textbf{67}, 042313.
%
% \bibitem{Cirelson1980}
%   Cirel'son B S (1980).
%   Quantum generalizations of Bell's inequality.
%   \textit{Letters in Mathematical Physics} \textbf{4}, 93--100.
%
% \bibitem{Childs2003}
%   Childs A M, Haselgrove H L, Nielsen M A (2003).
%   Lower bounds on the complexity of simulating quantum gates.
%   \textit{Physical Review A} \textbf{68}, 052311.

% ── End standalone ─────────────────────────────────────────────
% \end{document}
